% Conclusion
% Covers: summary of contributions, RQ answers, future work, final remarks
% Do not introduce new results here. Connect future work to limitations in Discussion.

\label{sec:conclusion}

\subsection{Summary of Contributions}

PROMPT introduces a phase-based course meta-model and its open-source implementation as a reusable platform for administering project-based courses across disciplines.
The meta-model addresses a gap unmet by existing \acrshort{lms} platforms and exercise-focused tools:
the explicit modeling of the course lifecycle as an ordered sequence of composable, independently operable phase instances.
The four default phase implementations -- Application, Interview, Team Allocation, and Assessment -- cover the most common administrative requirements of project-based courses, with the Assessment Phase providing a unique formative feedback infrastructure that formalizes self-evaluation, peer evaluation, and configurable instructor assessment criteria.

\subsection{Research Questions Answered}

\begin{itemize}
    \item \textbf{RQ1:} The phase-based meta-model captures the administrative structure of project-based courses more completely than existing tools, as demonstrated by the feature comparison and multi-semester deployment evidence at \university{}.
    \item \textbf{RQ2:} Instructors and \acrshortpl{ta} perceive PROMPT as reducing administrative overhead through workflow consolidation, with the Assessment Phase's formative feedback structure identified as the most valued capability.
    \item \textbf{RQ3:} The student \acrshort{sus} survey confirms \missing{high} perceived usability and strong appreciation of the Assessment Phase's transparency;
    multi-institution adoption provides preliminary evidence of the plugin architecture's generalizability.
\end{itemize}

\subsection{Future Work}

Future research directions include:
\begin{itemize}
    \item Cross-disciplinary adoption studies to provide empirical evidence for the meta-model's generalizability to project-based courses outside computing education.
    \item Phase configuration templates and semester-to-semester course cloning to reduce the setup overhead identified as the primary instructor pain point.
    \item Longitudinal studies on the impact of PROMPT's formative assessment support on student learning outcomes and reflection quality, representing the highest-priority research direction.
\end{itemize}

\subsection{Final Remarks}

PROMPT is available as open-source software at \url{https://github.com/prompt-edu/prompt} and is actively maintained by the \group{} research group at \university{}.
Institutions interested in deploying PROMPT for project-based courses -- in computing or any other discipline -- are invited to contribute phase modules to the growing community library.
